
\documentclass[letterpaper, 10 pt,conference]{IEEEconf}
\IEEEoverridecommandlockouts
\overrideIEEEmargins
% The preceding line is only needed to identify funding in the first footnote. If that is unneeded, please comment it out.
\usepackage{cite}
\usepackage{amsmath,amssymb,amsfonts}
\usepackage{algorithmic}
\usepackage{graphicx}
\usepackage{textcomp}
\usepackage{xcolor}
\usepackage{url}
%\def\BibTeX{{\rm B\kern-.05em{\sc i\kern-.025em b}\kern-.08em
%    T\kern-.1667em\lower.7ex\hbox{E}\kern-.125emX}}


\title{\LARGE \bf CADEMAS -- A Framework for Cooperative Automated Decision-Making Systems}

\author{
\IEEEauthorblockN{Pavel Novoa-Hernández}
\IEEEauthorblockA{\textit{Dept. Computer Engineering and Systems} \\
\textit{University of La Laguna}\\
La Laguna, Spain \\
pnovoahe@ull.edu.es}
\and
\IEEEauthorblockN{David A. Pelta, Mariia Godz, José L. Verdegay}
\IEEEauthorblockA{\textit{Dept. Computer Science and AI} \\
\textit{Universidad de Granada}\\
Granada, Spain \\
\{dpelta, mariiagodz, verdegay\}@ugr.es}
\and
\IEEEauthorblockN{Dionisio Buendia-Carrillo}
\IEEEauthorblockA{\textit{Dept. Financial Economics and Accountability} \\
\textit{Universidad de Granada}\\
Granada, Spain \\
dbuendia@ugr.es}
}



\author{Pavel Novoa-Hernández$^{1}$, David A. Pelta$^{2}$, Mariia Godz$^{2}$, José L. Verdegay$^{2}$, Dionisio Buendia-Carrillo$^3$
\thanks{$^{1}$P. Novoa is with Dept. Computer Engineering and Systems, University of La Laguna,La Laguna, Spain. 
{\tt\small pnovoahe@ull.edu.es}}
\thanks{$^{2}$D. Pelta, M. Godz, J.L. Verdegay are with Dept. Computer Science and AI, Universidad de Granada, Granada, Spain.
{\tt\small\{dpelta, mariiagodz, verdegay\}@ugr.es}}
\thanks{$^3$ D. Buendia-Carrillo is with Dept. Financial Economics and Accountability, Universidad de Granada, Granada, Spain.{\tt\small{dbuendia@ugr.es}}}
}


% \author{Author A$^{1}$, Author B$^{2}$, Author C$^{2}$, Author D$^{2}$, Author E$^3$
% \thanks{$^{1}$Author A is with Dept. xxxx, University of yyy. 
% {\tt\small homer@simpson.com}}
% \thanks{$^{2}$Author A, B, C are with Dept. zzzz, University of dddd, .
% {\tt\small\{marge,lisa,bart\}@simpson.com}}
% \thanks{$^3$ Author E is with Dept. qqq, University of ppp, .{\tt\small{who@simpson.com}}}
% }


\begin{document}
\maketitle

\begin{abstract}
The growing ubiquity of data and artificial intelligence has made it increasingly easy to construct decision models that automate or assist human decision-making. As organizations and infrastructures rely on multiple, heterogeneous \textit{Automated Decision-Making} (ADM) systems, a new challenge emerges: how to integrate and reconcile their outputs into coherent, explainable, and context-aware collective decisions. 
This paper proposes a new conceptual and architectural framework, the \emph{Cooperative Automated Decision-Making System} (CADEMAS), as a foundation for the study and design of cooperative decision architectures. 
CADEMAS formalizes the interaction among individual ADM systems through an Input Manager, an Output Manager, and a Context Model, allowing the transformation and aggregation of either homogeneous or heterogeneous decision outputs. 
CADEMAS generalizes cooperation across potentially heterogeneous decision agents, data types and goals, going beyond ensemble learning schemes, which combine homogeneous models over identical data.  
We argue that CADEMAS represents a necessary paradigm shift: from building isolated decision models to engineering cooperative ecosystems of decision-making entities. 
This paper outlines the theoretical basis, architectural elements, and implications for future research.
\end{abstract}

% \begin{IEEEkeywords}
% Automated Decision-Making (ADM), Cooperative Decision Systems, Decision Integration, Decision Architecture, Multi-Agent Systems, Artificial Intelligence.
% \end{IEEEkeywords}

\section{Introduction}
Automated decision making (ADM) is defined as ``a (computational) process, including AI techniques and approaches, that, fed by inputs and data received or collected from the environment, can generate, given a set of predefined objectives, outputs in a wide variety of forms (content, ratings, recommendations, decisions, predictions, etc)'' \cite{ELI_ADM_Principles_2022}.

Automated Decision-Making (ADM) systems  are now embedded in critical infrastructures and organizational processes, ranging from financial risk assessment and energy management to healthcare triage, transportation logistics, and public administration. The steady progress in machine learning, optimization, natural language processing, and computational statistics has made it increasingly simple to construct accurate, domain-specific automated decision models. As a result, modern decision ecosystems rarely rely on a single model; instead, they incorporate multiple ADM systems, often developed independently, trained on different data sources, and driven by distinct modeling assumptions or operational objectives.

Although constructing a single ADM system has become straightforward, combining their contributions into a coherent and reliable decision process remains a fundamental challenge. Real-world decision-making seldom arises from a single perspective. Decisions are influenced by heterogeneous signals, uncertain or incomplete information, cross-cutting criteria, and contextual, institutional, or regulatory constraints that evolve over time. In practice, organizations accumulate predictive models, diagnostic tools, rule-based systems, simulation components, and more recently, large language models (LLMs). Yet these systems typically operate in isolation, without a principled mechanism for cooperation, communication, or consensus. The absence of structured cooperation limits the robustness, transparency, and contextual alignment of the decisions produced by automated pipelines.

Existing approaches, like the ones shown later in Table \ref{tab:cade_example}, only partially address this situation. Traditional one-person decision-making problems assume a single rational agent optimizing a consistent objective function, an assumption unsuitable for settings where information, objectives, and decision processes are distributed. Multi-Agent Systems introduce cooperation, negotiation, and shared environments, but they generally rely on homogeneous representational structures and rarely incorporate institutional or regulatory contexts as explicit components of the decision process. Multi-Agent Reinforcement Learning enables emergent collective behavior, but cooperation derives from policy learning rather than from the combination of heterogeneous decision outputs, and it lacks mechanisms for integrating domain rules or external contextual evaluations.

Decision-support systems and multi-criteria decision-making methods provide structured ways to integrate diverse criteria, yet they typically assume a human-centric workflow and do not offer principles for the automated cooperation of independent computational agents. Ensemble learning, despite its demonstrated effectiveness, restricts itself to homogeneous model families with comparable input and output structures, focusing on improving accuracy rather than enabling cooperation among models with incompatible representations, goals, or semantics.

These limitations reveal four structural gaps in current approaches to automated decision-making. First, existing frameworks rarely consider the heterogeneity of data types, decision representations, and inference mechanisms that characterize modern ADM ecosystems. Second, they lack explicit mechanisms for semantic translation or communication among independently designed decision models. Third, most omit the role of external factors, such as regulatory frameworks, institutional priorities, ethical constraints, or situational factors, which often determine whether a technically sound decision can be executed. Finally, they do not incorporate iterative decision revision triggered by contextual misalignment, even though such iterations are common in human decision workflows.

As organizations increasingly rely on automated infrastructures, a framework is needed to define how heterogeneous ADM systems can cooperate, either indirectly - by producing independent outputs that are later harmonized - or directly, by exchanging information, refining each other’s assessments, or reconciling partial recommendations. Both ways of cooperation are necessary in real-world environments, and a general framework should be capable of accommodating them without imposing restrictive assumptions about model families, communication protocols, or decision representations.

The Cooperative Automated Decision-Making System (CADEMAS) framework proposed in this paper addresses these challenges by offering a unified conceptual and architectural view of cooperative automated decision-making. CADEMAS identifies three fundamental components that enable cooperation among ADM systems: a set of heterogeneous decision-making units, a semantic decision integration layer capable of translating, normalizing, and aggregating their outputs, and a contextual supervisor that evaluates the aggregated decision according to external constraints and triggers revisions when required. This structure formalizes how cooperation may occur and clarifies the role of context as a decisive element in modern automated decision systems.

The objective of this work is threefold. First, to motivate the need for cooperative, context-aware automated decision architectures and to analyze the limitations of existing approaches from this perspective. Second, to introduce CADEMAS as a general and flexible framework that captures the essential components of cooperation among ADM systems and third, to illustrate the practical relevance of the framework through a decision scenario involving the prediction of insolvency risk for Spanish small and medium-sized enterprises, a domain in which heterogeneous information sources and normative constraints coexist.

The remainder of this paper is structured as follows. Section \ref{sect:background} presents the notions of decision, decision-making agent, and Automated Decision-Making system as used in this work. Section \ref{sect:framework} introduces the CADEMAS framework and describes its architectural components. Section \ref{sect:examples} provides an overview of existing systems that can be fitted under the CADEMAS framework. 
Then, section \ref{sect:smes} presents an illustrative case study situated within the context of SME insolvency prediction. Finally, section \ref{sect:conclus} discusses implications, potential applications, and limitations. 

\section{Basic Definitions\label{sect:background}}

CADEMAS are designed to provide a unified view of how heterogeneous decision-making entities can cooperate to produce contextually aligned and semantically coherent decisions. Before introducing the formal architecture, we clarify the notions of decision, decision-making agent, and Automated Decision-Making system as used in this work.

\subsection{Decisions and Decision-Making Agents}

From a theoretical point of view \cite{LamataPeltaVerdegay2020}, and considering a single decision maker, the definition of a decision problem requires the following essential elements:
$$(X, I, E,C,\leq, K)$$
% \begin{itemize}
%     \item  One decision-maker, either an individual or a group,
%     \item a set of actions on which the decision maker can choose,
%     \item a set, called the Environment, which is constituted by the situations (states) that the decision maker can encounter when choosing, and cannot control,
%     \item a set of consequences, associated with each action and each state,
%     \item a criterion that orders the consequences,
%     \item the nature of the information available,
%     \item the (social) framework in which the decision-making process takes place.
% \end{itemize}

%Thus, if we assume in all that follows an individual decision maker, a decision problem is described by a sextet $$(X, I, E,C,\leq, K)$$ 
\noindent
where $X$ is the set of possible actions for the decision maker; $I$ is the available information; the environment $E$ which is constituted by the situations (states) that the decision maker can encounter when choosing, and cannot control; the set $C$ of the consequences of actions on every state; the criterion $\leq$ that orders the consequences; and the framework $K$, which is the context in which the decision maker decides.

%cambié $\mathcal{A}$ por $\mathcal{S}$. la primera la usabamos después para otra cosa.
A \textit{decision} is understood here as a commitment to a course of action or to the selection of one option among several admissible alternatives, under a given representation of the environment, constraints, and goals. Formally, given a set of admissible actions $\mathcal{S}$, a decision is an element $d \in \mathcal{S}$ or a structured object encoding a preference, ranking, probability distribution, or policy over $\mathcal{S}$. 
It should be remarked that a decision does not require the execution of an action; it is the outcome of a deliberation process that \emph{could} lead to an action.

A \textit{decision-making agent} is any entity capable of making decisions according to an internal model or set of rules. Agents may be humans, algorithms, institutions, or hybrid combinations thereof. In this work, we focus on ADM systems while allowing human agents to act at supervisory or contextual levels if necessary.

\subsection{Automated Decision-Making Systems}

As stated before, an ADM system is a software system that runs computational processes which, given some input representation of a situation, produces a decision or recommendation with minimal or no human intervention. 


% Formally, an ADM is a function
% \[
%     ADM : IT \rightarrow OT,
% \]
% where $IT$ and $OT$ denote admissible input and output types

The input/output information may include numerical vectors, symbolic rules, images, fuzzy sets, natural language, or probability distributions.

ADM systems vary widely in design and purpose. Examples include machine learning classifiers, regression models, optimization algorithms, rule-based systems, simulation engines, instance-based reasoning systems, or large language models acting as decision-generating agents. Importantly, ADM systems generally operate isolated; they do not communicate their internal states, nor do they coordinate with other ADM systems unless explicitly engineered to do so.

\section{CADEMAS: A Cooperative Architecture for Automated Decision-Making\label{sect:framework}}

CADEMAS provides a conceptual and architectural framework for structuring cooperation among, in principle, heterogeneous Automated Decision-Making units.

Rather than prescribing a particular learning method, rule formalism, or 
coordination protocol, CADEMAS specifies how diverse decision-producing 
components can be harmonized into a unified and context-aware final decision. 
The framework is organized around three main layers: ADM Units, the Decision 
Integration Layer, and the Contextual Modulator; each with clearly delineated 
roles.

A central premise of CADEMAS is that an \emph{ADM} system is not necessarily a 
full-fledged agent capable of making an operational decision. Instead, an ADM system can be an automated component that produces a \emph{partial decision signal}: a score, prediction, rule activation, explanation, or suggested action. These
signals may lack semantic homogeneity and are not considered final decisions
until processed by the integration and contextual layers. This distinction 
clarifies that the \emph{decisor} in CADEMAS is the system as a whole.

Formally, a CADEMAS instance is defined as the tuple
\[
    \mathcal{S} = \langle \mathcal{IM}, \mathcal{A}, \mathcal{OM}, \mathcal{C}, \mathcal{T} \rangle,
\]
where $\mathcal{IM}$ is the Input Manager, $\mathcal{A}$ is the set of ADMs,
$\mathcal{OM}$ the Decision Integration Layer, $\mathcal{C}$ the Contextual
Modulator, and $\mathcal{T}$ the universe of admissible data types and 
representations. 

%Figure~\ref{fig:cademas_architecture} provides an overview of the main components of CADEMAS, which are described in detail in the following sections.

% \begin{figure}[t]
%     \centering
%     \includegraphics[width=0.8\columnwidth]{cademas_architecture.pdf}
%     \caption{CADEMAS architecture showing the interaction between ADM units, 
%     the integration layer, and the contextual modulator.}
%     \label{fig:cademas_architecture}
% \end{figure}

\subsection{Input Manager}

Although CADEMAS is not tied to specific data sources, it is important to 
distinguish the two conceptually different forms of information entering the 
architecture. The \emph{Input Manager} $\mathcal{IM}$ is responsible for 
routing, transforming and validating both:

\begin{enumerate}
    \item \textit{Training or calibration data}, used to construct or update the 
    internal models of each $ADM_i$. These inputs are consumed offline and may 
    include historical financial records, domain annotations, or expert-labelled 
    cases. They do not participate directly in a concrete decision.

    \item \textit{Decision-case data}, provided whenever a new decision request 
    is evaluated. These include the observable attributes of the entity under 
    analysis (e.g., a firm) and any contextual variables required by the 
    modulator. Unlike training data, these inputs are routed to the ADMs in 
    inference mode and to the contextual component simultaneously.
\end{enumerate}

CADEMAS therefore separates \emph{model construction} from 
\emph{decision operation}: a system may be periodically recalibrated using 
large training corpora, while individual decision cases may be evaluated with 
minimal, contextually relevant information. This distinction reinforces the 
role of $\mathcal{IM}$ as a unifying gateway through which heterogeneous data 
types enter the architecture in a coherent manner.

\subsection{ADM systems layer}

The set of ADM systems is denoted:
\[
    \mathcal{A} = \{ADM_1, \dots, ADM_n\}.
\]


Each $ADM_i$ receives a data input representation $d_i \in \mathcal{T}$ generated by the Input Manager and produces an output $o_i \in \mathcal{T}$. For example, an $ADM_i$ system may correspond to predictive models generating probabilities or class scores, optimization modules returning candidate solutions, simulation components providing scenario estimates, rule-based systems encoding normative or domain rules, LLM-based units doing symbolic or textual reasoning, to cite just a few.


Importantly, CADEMAS does \emph{not} assume that $ADM_i$ themselves must make a final decision. Their role is to contribute evidence, judgments, or 
preliminary recommendations. Cooperation among ADM systems may therefore be 
\emph{indirect}, with each operating independently and contributing its 
own signal, or \emph{direct}, with ADM systems sharing intermediate representations
or refining each other's outputs. CADEMAS allows both modes without imposing
any specific communication protocol.

\subsection{Decision Integration Layer}

Given the heterogeneity of the $ADM_i$ outputs 
\[
    \mathcal{O} = \{o_1, \dots, o_n\},
\]
a dedicated integration mechanism is required to produce a unified decision 
object $D$. The Decision Integration Layer is formalized as a mapping:
\[
    \Phi : \mathcal{O} \rightarrow D,
\]
and it performs two essential functions.

\paragraph*{Semantic Harmonization}
ADM outputs may differ in type and semantics (e.g., probabilities, scores, 
rule activations, text). The integration layer transforms each $o_i$ into 
a shared representational ontology, which may include normalized scores, 
linguistic categories, uncertainty intervals, or structured decision tokens. 
In settings involving textual or symbolic information, LLMs can facilitate 
semantic harmonization.

\paragraph*{Aggregation and Synthesis}
Once harmonized, the signals are aggregated into a decision object $D$. Possible mechanisms include weighted voting, multi-criteria decision making, consensus reaching, belief merging, or LLM-based meta-reasoning. The output $D$ is a \emph{raw} decision proposal: coherent, syntactically unified, but not yet context-adjusted.

\subsection{Contextual Modulator}

The Contextual Modulator distinguishes CADEMAS from classical 
multi-agent or ensemble-based architectures. Unlike a validator that 
accepts or rejects the integrated decision, the contextual component 
\emph{modulates} the raw decision according to exogenous constraints, norms,
or situational factors. 

As stated in \cite{Dourish2004}, the context is a form of information: it is delineable and stable, and there is a separation between the context and the problem being solved (the problem is addressed within a context).
    
The authors in \cite{PerezCanedo2022} proposed modeling context in optimization problems using fuzzy logic-based propositions (rules). In this way, the degree to which the solution satisfies the context rules can be measured. 

A context $\mathcal{C}$ may encode, for example, legal or regulatory requirements, organizational policies or fairness criteria, budgetary or resource limits, macroeconomic or environmental conditions, and so on.


Formally, the modulator is a mapping:
\[
    M : D \times \mathcal{C} \rightarrow D,
\]
where $M$ adjusts, reinterprets, or conditions the raw decision to produce 
a context-aware final decision:
\[
    D^{\text{final}} = M(D^{\text{raw}}, \mathcal{C}).
\]

The modulation may involve altering risk categories, redefining the type or intensity of the recommended action, attaching conditions, or modifying its timing. This approach avoids rigid accept/reject dynamics and better reflects how decisions are made in real-world institutions, where contextual and policy-based nuances regularly reshape technically derived recommendations.

Possible realizations of $M$ include rule-based modulating layers,  constraint-aware adjustment operators, policy-driven transformations, or LLM-based modulators capable of interpreting textual policy documents and producing contextually consistent refinements.

In summary, CADEMAS conceptualizes decision-making as a cooperative process where heterogeneous ADM systems provide partial decision signals, a unified decision is synthesized through semantic and algorithmic integration, and the final action emerges through context-sensitive modulation rather than binary validation. This layered structure enables both flexibility and accountability in complex decision environments.


\section{Examples of CADEMAS-like systems\label{sect:examples}}

Although the label \emph{CADEMAS} is new, many research efforts implicitly develop cooperative decision-making mechanisms that naturally fit within the structure of our framework. These works emerge from diverse communities such as multi-agent reinforcement learning, distributed constraint optimization, LLM-based orchestration, and in human-AI collaborative governance, and all of them depart from a common intuition: complex decisions often require the interaction of several autonomous decision entities rather than reliance on a single model or policy. What these proposals lack, however, is a unifying abstraction that separates domain-specific techniques from the more general structural pattern they instantiate. CADEMAS provides such an abstraction by identifying the architectural invariants that underlie these heterogeneous systems.

Table~\ref{tab:cademas_mapping} summarizes some recent works, showing how diverse proposals define the components of CADEMAS through different representational choices and operational objectives. The purpose is not to impose rigid boundaries or reclassify existing systems but to highlight the structural commonality across them. In all cases, independent decision entities operate over heterogeneous representations, produce partial outputs that must be harmonized through aggregation mechanisms, and are subject to contextual constraints that may confirm or invalidate the integrated decision.

CADEMAS contributes value by making this underlying structure explicit and model independent. It clarifies the role of representational management (the input manager), decision autonomy (the ADM systems layer), aggregation and semantic alignment (the integration layer), and external normative or situational modulation (the contextual supervisor). This articulation provides not only a unifying perspective on existing cooperative decision systems but also a principled foundation for designing new ones in domains where heterogeneous models, multi-source information, and context-sensitive constraints coexist.

\begin{table*}[h!]
\centering
\renewcommand{\arraystretch}{1.3}
\caption{A showcase of CADEMAS-like examples in the literature. The main components are briefly described}
\label{tab:cademas_mapping}
\begin{tabular}{|p{3.2cm}|p{3.2cm}|p{3.2cm}|p{3.2cm}|p{3.2cm}|}
\hline
\textbf{Work} & \textbf{Input Manager $\mathcal{IM}$} & \textbf{Agents $\mathcal{A}$} & \textbf{Output Manager $\mathcal{OM}$} & \textbf{Context $\mathcal{C}$} \\
\hline
Mingyue (2025) \cite{Cai2025} & Prompt templates, role conditioning & LLM planner, verifier, critic & Iterative critique-consensus & Scenario goals, constraints \\
\hline
CoDrivingLLM (2024) \cite{CoDrivingLLM} & Sensory-to-text transformations & LLM modules (risk, rules, planning) & Rule-based safety consensus & Driving environment, safety limits \\
\hline
Trust-MARL (2025) \cite{Trust-MARL2025} & State extraction, trust features & MARL agents & Trust-weighted action fusion & Cooperative temporal objectives \\
\hline
DITURIA (2022) \cite{dituria2022}& Mapping mixed preference data & DCOP agents & Preference aggregation & Multi-criteria decision space \\
\hline
HEC (2025) \cite{HECframework} & Ethical claim processing & Norm-evaluating agents & Ledger-based consensus & Ethical principles and norms \\
\hline
\end{tabular}
\end{table*}

\section{A CADEMAS for small and medium-sized enterprises rescue prioritization\label{sect:smes}}

To illustrate the operation of CADEMAS, we consider a case study based on a real dataset of Spanish small and medium-sized enterprises (SMEs) \cite{NavarroGalera2025}. The motivating problem is the early prioritization of public intervention for firms at risk of insolvency \cite{NavarroGalera2025}. Public authorities must allocate limited rescue resources under uncertainty, balancing technical risk estimates with sectorial relevance, socio-economic impact, and contextual feasibility \cite{PerezCanedo2022}. CADEMAS provides a structured framework for this multimodel, context-aware prioritization task.

The dataset, obtained from the Spanish Mercantile Registry (\url{https://www.rmc.es/}), includes longitudinal financial and non-financial information for 10,000 SMEs (5,000 insolvent and 5,000 solvent), with up to twenty annual observations per firm. Missing values were handled through interpolation techniques and replacement with central tendency measures. Model performance was evaluated on a hold-out set comprising 25\% of the cases. For illustration, five representative firms (A–E) are intentionally selected from the 25\% hold-out test set to demonstrate the CADEMAS prioritization process. Each firm is described by fourteen attributes: four non-financial variables (e.g., sector, audit opinion, size, legal form) and ten financial ratios.


\subsection{ADM systems Layer}

%CADEMAS defines the ADM layer as a collection of heterogeneous decision-making units capable of producing partial assessments of the same problem.
As stated before, each ADM system in this layer can be based on different modeling paradigms depending on the data and the task: for instance, classical time-series models, gradient boosting, recurrent neural networks, survival analysis, or dynamic panel econometrics. Here, we select four representative ADM systems to highlight complementary perspectives and how their outputs can be integrated within CADEMAS.\\

\noindent
\textit{ADM$_1$: Gradient Boosting on Time-Series Features}, summarizes short-term financial micro-series through statistics such as mean, slope, volatility, and autocorrelation, and produces a probability of insolvency at one year. This system demonstrates how feature-based supervised learning can capture temporal patterns without modeling full sequences. Performance measure: AUC-ROC.\\

\noindent
\textit{ADM$_2$: LSTM Network}, processes the raw annual sequences of liquidity, leverage, profitability, and payment-delay histories, providing as output the probability of \textit{Insolvent}/\textit{Solvent}. This illustrates a sequence-based deep learning approach capable of capturing long-term dependencies that simpler feature-based models may miss. Performance measure: F1-score.\\

\noindent
\textit{ADM$_3$: Survival Analysis}, a Cox proportional hazards model estimates one-, three-, and five-year hazard rates. Here we focus on the one-year marginal hazard to generate a comparable risk score. Survival models offer a time-to-event perspective, complementing probability-focused classifiers. Performance measure: Concordance Index (C-index).\\

\noindent
\textit{ADM$_4$: Econometric Dynamic Panel Model}, a fixed-effects logistic panel regression captures firm-level heterogeneity and temporal dependencies using lagged financial ratios. This ADM system emphasizes interpretability and longitudinal modeling, showing how econometric methods can fit within CADEMAS. Performance measure: Accuracy.


%These four ADM systems provide a variety of decision signals that are fed to the integration layer.

\subsection{Decision Integration Layer}

The decision integration layer will combine the heterogeneous outputs of ADM systems into a single, actionable signal. In general, integration could be achieved through weighted sums, rank-based aggregation, consensus voting, etc. As a first approach, we compute a risk $R_i$ for a company $i$ using a weighted sum of normalized risk scores, where the weights are derived from the predictive performance measures (defined in $[0,1]$) associated with each ADM system:
\[
    R_i = \sum_{j=1}^{4} w_j \, o_{ij},
\]
\noindent
The weights \(w_j\) are proportional to the corresponding validation performance indicators (e.g., AUC-ROC, F1-score, C-index, Accuracy) and are normalized so that their sum equals 1. In our example, we obtain $w = (0.429, 0.353, 0.131, 0.087)$.

\subsection{Contextual Modulator: Public-Interest Rescue Criteria}

CADEMAS treats context as an external, modulated input. The framework itself does not impose a specific formalism; context can be modeled through crisp rules, utility functions, or fuzzy logic, among others. Following recent work \cite{PerezCanedo2022}, fuzzy logic provides a convenient way to quantify degrees of relevance across multiple criteria.

In our case, the context represents public-interest criteria for potential SME rescue. We define three fuzzy dimensions: employment impact, strategic sector relevance, and regional/social vulnerability. 
Then, the degree to which a firm satisfies each dimension/criterion (i.e. a membership value \(\mu \in [0,1]\) is calculated.

To aggregate these values into a single contextual relevance score $C_i$, we apply a conservative fuzzy \textit{AND} (minimum operator). In all cases, the membership functions are defined as triangular functions whose parameters are derived directly from the empirical ranges and semantic definitions of the corresponding attributes in the dataset:
\begin{equation*}
    C_i = \min\{\mu_\text{emp}(i), \mu_\text{sec}(i), \mu_\text{soc}(i)\}.
\end{equation*}
\noindent
This ensures that a firm is considered contextually relevant only if it performs adequately across all dimensions. Other aggregation methods could be used, such as weighted fuzzy sums or threshold-based rules, depending on the intended policy interpretation.


\subsection{Final CADEMAS Decision: Combined Priority Index}

\begin{table*}[tb]
\centering

\caption{Summary of CADEMAS output for the analyzed SMEs.}
\renewcommand{\arraystretch}{1.3}
\label{tab:cade_example}
\begin{tabular}{|c|c|c|c|c|c||c|c|c|c||c|c|}
\hline
\textbf{Firm} & \textbf{ADM$_1$} & \textbf{ADM$_2$} & \textbf{ADM$_3$} & \textbf{ADM$_4$} & \textbf{Weighted Risk} & 
\textbf{Employ.} & \textbf{Sector} & \textbf{Social} & \textbf{Context} & 
\textbf{Priority Index} & \textbf{Rank} \\
 & TS-GBM & LSTM & Survival & Econometric & \(R\) & \(\mu_{emp}\) & \(\mu_{sec}\) & \(\mu_{soc}\) & \(C\) & \(P\) & \\
\hline
A & 0.978 & 0.870 & 0.671 & 0.513 & 0.859 & 0.900 & 0.800 & 1.000 & 0.800 & 0.830 & 2 \\
B & 0.897 & 0.771 & 0.631 & 0.443 & 0.778 & 0.600 & 0.400 & 0.500 & 0.400 & 0.589 & 4 \\
C & 0.912 & 0.671 & 0.712 & 0.688 & 0.781 & 1.000 & 0.900 & 0.900 & 0.900 & 0.841 & 1 \\
D & 0.951 & 0.820 & 0.119 & 0.432 & 0.751 & 0.400 & 0.600 & 0.700 & 0.400 & 0.575 & 5 \\
E & 0.981 & 0.721 & 0.213 & 0.541 & 0.750 & 0.800 & 1.000 & 0.600 & 0.600 & 0.675 & 3 \\
\hline
\end{tabular}
\end{table*}

Finally, we combine the ADM-derived risk \(R_i\) with the contextual relevance \(C_i\) to produce a prioritization index:
% \[
%     P_i = \frac{1}{2} \left( R_i + C_i \right),
% \]
\[
    P_i = \alpha \times R_i + (1-\alpha) \times C_i,
\]
where both components are normalized to \([0,1]\). Here, $\alpha = 0.5$ so the index \(P_i\) captures firms that are simultaneously at risk and are strategically relevant for public intervention. Higher values indicate higher priority. This formulation is just one way to integrate the final decision. Again, the choice of ADM systems, integration function, or contextual modeling could vary, but the overarching CADEMAS architecture remains applicable.





Table~\ref{tab:cade_example} summarizes the outputs produced by CADEMAS. Each firm (rows) is evaluated by the four ADM systems, which provide heterogeneous insolvency-risk signals (columns $ADM_i$). As expected, the models do not exhibit perfect agreement: $ADM_1$ tends to assign higher predicted risks, while the remaining ADM systems show more moderate values and differ in their sensitivity to the underlying financial patterns. These signals are aggregated through performance-based weights into the \textit{Weighted Risk} column, in which Firm~\textit{A} emerges as the firm with the highest technical insolvency risk.

Contextual criteria (e.g., employment impact, Sector relevance, and Social vulnerability) are represented by their fuzzy membership degrees (columns \textit{Employ.}, \textit{Sector}, and \textit{Social}, Table \ref{tab:cade_example}). These values reflect the extent to which each firm satisfies the public-interest criteria defined earlier. Firm~C stands out as the one achieving the highest joint contextual relevance, as reflected in the \textit{Context} column (computed using a conservative fuzzy \textit{AND}).

Finally, the \textit{Priority Index} and its derived ranking (last two columns) illustrate how CADEMAS integrates technical risk with public-interest relevance. Notably, the resulting order differs from the ranking based solely on aggregate risk: Firm~C becomes the top priority for intervention, displacing Firm~A despite the latter having the highest risk of insolvency. This demonstrates the essential role of the contextual modulator in preventing purely technical assessments from dominating public-sector decision processes.

\section{Discussion and Implications\label{sect:conclus}}
We proposed here a new conceptual and architectural framework, the \emph{Cooperative Automated Decision-Making System} (CADEMAS), as a foundation for the study and design of cooperative decision architectures. 

The framework integrates many already existing tools that are available or studied in isolated terms, like multi-model tools, aggregation mechanism, context-aware evaluations, ensemble-like methods, etc. 

CADEMAS formalized the interaction among individual ADM systems through an Input Manager, an Output Manager, and a Context Model, allowing the transformation and aggregation of homogeneous or heterogeneous decision outputs. 

The application example showed just a few of the design alternatives that can be used under the CADEMAS scheme. In particular, the ways in which ADM systems output can be combined and how the context can be modeled, open the way to enhance the automated decision-making process, providing more informed decisions and better suited to the context at hand. 

The use of several ADM systems allows for at least two working modes. The first was shown in the example: each ADM system solved the problem from a different perspective and using different kinds of information. The second one arises when there is more than one ADM system available to solve the problem. Instead of searching for the ``best'' ADM, CADEMAS provides a framework to use all of them in order to provide more solid and informative decisions as they are supported by several individual ``decisors''.


Overall, the theoretical basis, architectural elements, and implications for future research on CADEMAS have been established.
%This study laid the theoretical and architectural elements for CADEMAS while identifying key areas for future investigation.

\section*{Acknowledgment}
Mariia Godz has a pre-doctoral research contract associated with the project “Study, Analysis and Evaluation of Cooperative Automated Decision-Making Systems (CADEMAS)”, PID2023-146575NB-I00, funded by MCIU/AEI/I0.13039/501100011033 and by FSE+. The authors acknowledge support of the previous project.


%\bibliographystyle{abbrv}
%\bibliography{biblio}
\begin{thebibliography}{10}

\bibitem{Cai2025}
X.~Cai, et.al.
\newblock MINGYUE: A multi-agent system for human-ai collaborative decision-making under safety constraints in complex power grids.
\newblock In {\em 2025 5th Power System and Green Energy Conference (PSGEC)}, pages 434--439, 2025.

\bibitem{Dourish2004}
P.~Dourish.
\newblock {What we talk about when we talk about context}.
\newblock {\em Personal and Ubiquitous Computing}, 8(1):19--30, 2004.

\bibitem{ELI_ADM_Principles_2022}
{ELI Innovation Paper}.
\newblock {Guiding Principles for Automated Decision-Making in the EU}.
\newblock {European Law Institute (ELI)}, 2022.

\bibitem{CoDrivingLLM}
S.~Fang, J.~Liu, M.~Ding, Y.~Cui, C.~Lv, P.~Hang, and J.~Sun.
\newblock Towards interactive and learnable cooperative driving automation: A large language model-driven decision-making framework.
\newblock {\em IEEE Transactions on Vehicular Technology}, 74(8):11894--11905, 2025.

\bibitem{dituria2022}
H.~Ibro, G.~Mahala, S.~Pulawski, S.~Harvey, A.~A. Miller, A.~Ghose, and H.~K. Dam.
\newblock Dituria: A framework for decision coordination among multiple agents.
\newblock In {Proceedings of Principles and Practice of Multi-Agent Systems (PRIMA):}, page 492–508. Springer-Verlag, 2022.

%\newblock In {\em PRIMA 2022: Principles and Practice of Multi-Agent Systems: 24th International Conference Proceedings}, page 492–508. Springer-Verlag, 2022.

\bibitem{HECframework}
R.~K. Lakshmana~Perumal.
\newblock Holistic ethical commons: A dynamic framework for ethical decision making in multi-agent ai systems.
\newblock In {\em 5th Intelligent Cybersecurity Conference (ICSC)}, pp 378--384, 2025.

\bibitem{LamataPeltaVerdegay2020}
M.~T. Lamata, D.~A. Pelta, and J.~L. Verdegay.
\newblock The role of the context in decision and optimization problems.
\newblock In {\em Fuzzy Approaches for Soft Computing and Approximate Reasoning: Theories and Applications}, pages 75--84. Springer International Publishing, 2020.

\bibitem{NavarroGalera2025}
A.~Navarro-Galera, J.~Lara-Rubio, P.~Novoa-Hernández, and C.~A. Cruz~Corona.
\newblock Using decision trees to predict insolvency in spanish smes: Is early warning possible?
\newblock {\em Computational Economics}, 65(1):91–116, Mar. 2025.

\bibitem{Trust-MARL2025}
J.~Pan, T.~Wang, C.~Claudel, and J.~Shi.
\newblock Trust-marl: Trust-based multi-agent reinforcement learning framework for cooperative on-ramp merging control in heterogeneous traffic flow, 2025.

\bibitem{PerezCanedo2022}
B.~P{\'{e}}rez-Ca{\~{n}}edo, C.~Porras, D.~A. Pelta, and J.~L. Verdegay.
\newblock Modeling contexts as fuzzy propositions in optimization problems.
\newblock {\em {IEEE} Transactions on Fuzzy Systems}, 31(5):1474--1483, may 2023.

\end{thebibliography}

\end{document}


%%%%% AQUI ESTÁ EL CONTENIDO ORIGINAL COMPLETO DE ESTA SECCION
\section{Positioning CADEMAS in the Current Literature}

Although the label \emph{CADEMAS} is new, many research efforts implicitly develop cooperative decision-making mechanisms that naturally fit within the structure of our framework. These works emerge from diverse communities—multi-agent reinforcement learning, distributed constraint optimization, LLM-based orchestration, and human–AI collaborative governance—but they share a common intuition: complex decisions often require the interaction of several autonomous decision entities rather than reliance on a single model or policy. What these proposals lack, however, is a unifying abstraction that separates domain-specific techniques from the more general structural pattern they instantiate. CADEMAS provides such an abstraction by identifying the architectural invariants that underlie these heterogeneous systems.

A review of representative works illustrates this point. Multi-agent LLM systems for planning, such as those proposed by Mingyue et al. (2025), coordinate multiple specialized language models—typically a planner, critic, verifier, or simulator—whose interaction is mediated by prompt templates and iterative reasoning loops. Although the focus is on procedural task-solving, the underlying structure is clearly compatible with CADEMAS: multiple decision-making units producing partial recommendations, an integration mechanism synthesizing critiques and revisions, and contextual task goals constraining the final output.

A similar pattern appears in CoDrivingLLM (2024), a multi-agent architecture for autonomous driving in which several LLMs reason about perception, trajectory risk, and traffic rules. Sensory data are transformed into textual descriptions, processed by domain-specific LLM agents, and consolidated through safety-driven consensus rules. The context layer is here embodied by real-time environmental dynamics and legally codified driving constraints—again mirroring the CADEMAS decomposition.

In multi-agent reinforcement learning, the Trust-MARL framework (2025) introduces trust-aware cooperative policies in dynamic environments. Despite relying on continuous interaction and learning rather than on static decision proposals, the same logical structure persists: agents derive partial utilities or recommended actions, a trust-weighted mechanism aggregates them, and temporal or cooperative objectives act as the contextual evaluator guiding policy adjustment.

Distributed optimization systems, such as DITURIA (2022), follow a comparable architecture. Agents encode constraints or preferences and cooperate through distributed solving and preference fusion. A global feasibility or multi-criteria evaluator operates as a contextual supervisor that validates or rejects candidate solutions, effectively creating the same iterative loop that CADEMAS formalizes.

Other lines of work reveal analogous structures under different semantic interpretations. The HEC framework (2025), oriented toward ethical reasoning and collective norm revision, processes ethical claims through heterogeneous reasoning agents and consolidates them via ledger-based consensus, while societal norms or ethical principles act as contextual constraints. Similarly, in cognitive architectures such as MAIBL (2023), agents base decisions on instance retrieval and memory similarity, and combined outputs are evaluated in relation to situational compatibility or task requirements.

Table~\ref{tab:cademas_mapping} summarizes these correspondences, showing how diverse proposals instantiate the components of CADEMAS through different representational choices and operational objectives. The purpose is not to impose rigid boundaries or reclassify existing systems but to highlight the structural commonality across them. In all cases, independent decision entities operate over heterogeneous representations, produce partial outputs that must be harmonized through explicit or implicit aggregation mechanisms, and are subject to contextual constraints that may confirm or invalidate the integrated decision.

CADEMAS contributes value by making this underlying structure explicit and model-agnostic. It clarifies the role of representational management (the input manager), decision autonomy (the ADM layer), aggregation and semantic alignment (the integration layer), and external normative or situational modulation (the contextual supervisor). This articulation provides not only a unifying perspective on existing cooperative decision systems but also a principled foundation for designing new ones in domains where heterogeneous models, multi-source information, and context-sensitive constraints coexist.
